\section{Introduction}\label{sec:intro}

Polynomial optimization offers a principled route to globally informed robot
trajectory optimization. Once dynamics, contact complementarity and kinematic
constraints are written, or lifted, as polynomial equalities and inequalities,
the moment--sum-of-squares (SOS) hierarchy~\cite{lasserre2001global,parrilo2003semidefinite}
produces semidefinite programs (SDPs) whose optimal values are certified lower
bounds on the trajectory cost. Their solutions often lead to near-optimal
trajectories by extraction or rounding. Relaxations that exploit correlative and
term sparsity~\cite{waki2006sums,lasserre2006convergent,wang2021tssos,wang2022cstssos}
keep long horizons representable. SPOT~\cite{kang2025spot} has demonstrated this
for contact-rich manipulation, related relaxations are tight on several
contact-rich problems~\cite{graesdal2024tight}, and STROM~\cite{kang2024strom}
solves sparse second-order relaxations of polynomial trajectory problems to
certified suboptimality with a first-order (ADMM) SDP solver on the GPU.

Being able to write down such a relaxation does not make it easy to solve.
A sparse trajectory relaxation contains tens of thousands of small
positive-semidefinite (PSD) blocks and hundreds of thousands of equality
constraints; one of our PushT instances with horizon $N=100$ has $259{,}499$
equality rows and $63{,}885$ PSD blocks of order at most $13$. Interior-point
methods (IPMs) remain the method of choice when high accuracy is required, and
most of their time goes into forming and factoring the Schur-complement (Newton)
matrix on the equality constraints. Small PSD blocks do not by themselves make
this matrix cheap to factor. Coefficient-matching rows couple several blocks,
monomials in globally shared variables can span the whole horizon, and an
ordering chosen without knowledge of the model can create large dense fronts.
General-purpose sparse direct solvers on CPUs (as in MOSEK~\cite{mosek}) or GPUs
(as in cuDSS~\cite{cudss}) recover part of this structure through fill-reducing
orderings. They do not know which couplings are temporal, and they cannot
guarantee that a hierarchy remains compatible after the model-specific
simplifications that make elimination cheap. First-order GPU methods such as
STROM's ADMM avoid the Newton system altogether and scale well, but need many
iterations to reach high accuracy. We keep the second-order IPM and make its
Newton system cheap instead.

This paper connects the polynomial model to the graph that the linear solver
actually sees, and uses that connection to build an exact, certified,
GPU-friendly elimination scheme for the IPM Newton system. Our contributions are
as follows.

\begin{enumerate}[leftmargin=*,itemsep=2pt]
\item \textbf{Temporal separators from coefficient matching.} We define the
  \emph{resource graph} of an SOS relaxation, in which every PSD block or free
  variable induces a clique on the coefficient rows it touches. This graph
  contains the sparsity pattern of the Schur matrix and all of its symbolic fill
  (\cref{lem:envelope}). A temporal cut has a separator of size at most
  $\binom{s+2r}{2r}$, where $s$ is the number of polynomial variables shared
  across the cut and $r$ is the relaxation order (\cref{thm:monomial}). Under
  bounded temporal support, an explicit interval ordering gives a path
  decomposition of bounded width, and hence linear factorization work in~$N$
  (\cref{thm:interval}).
\item \textbf{A certified hierarchy with exact condensation.} We give an
  adjacent-interface certificate, checkable in time linear in the size of the
  Schur graph, under which independent elimination of the interior components
  yields an exact block-tridiagonal system with a global border
  (\cref{thm:chain}). The certificate validates aggressive separator pruning and
  temporal coarsening, which are justified algebraically by associativity of
  Schur complements (\cref{prop:assoc}). A compatible hierarchy always exists
  under the modeling assumptions (\cref{thm:spaced}).
\item \textbf{Complexity bounds that depend on the actual fronts.} We bound
  factorization work, storage and solve cost by accumulated local elimination
  fronts, component boundary sizes and retained interface widths
  (\cref{thm:main}). Under uniform local bounds this gives $O(N)$ numerical work
  and storage (\cref{cor:linear}) and $O(\log N)$ dependency depth
  (\cref{cor:depth}). We refine the bounds for unequal interfaces and for the
  supernodal plan that the GPU executes. We also state what the bounds do not
  cover: symbolic setup, the number of outer IPM iterations, and floating-point
  accuracy.
\item \textbf{A GPU interior-point solver.} We present MIST, an
  interior-point solver for SOS relaxations of multistage polynomial
  optimization, and describe its structured backend. It combines sparse-coordinate Schur assembly, certified symbolic
  planning, batched native supernodal partial factorization of repeated stage
  patterns, and a reduced solve by dense Cholesky, cyclic reduction or sparse
  factorization. Every direction is checked against the full operator, and the
  backend falls back to a general factorization when a check fails.
\item \textbf{Experiments at matched and strict accuracy.} On SPOT planning
  SDPs, the hierarchical backend is $1.19$--$1.22\times$ faster than full-system
  cuDSS on the same A100 at $N=30$ and $100$, and the reduced solve takes about
  $1\%$ of Newton time. On PushBox with $N$ up to $80$, MIST is
  $1.5$--$3.2\times$ faster than MOSEK~11 and BRISK~\cite{brisk} for $N\ge20$.
  For $N\ge5$, it is the only one of the three solvers whose solutions pass a
  strict $10^{-6}$ check on the original SDP. On a three-dimensional pendulum
  problem with $N=20$, MIST on one H100 converges $12.3\times$ faster than
  8-thread MOSEK.
\end{enumerate}

The graph-theoretic and block-elimination identities we use are
classical~\cite{rose1976algorithmic,george1973nested,vandenberghe2015chordal},
and linear per-iteration complexity for suitably structured sparse SDPs is known
for other formulations~\cite{zhang2020clique,zhang2024chordal}. We do not claim
the first linear-complexity sparse SDP method. Our contribution is the
specialization of these tools to SOS coefficient incidence, a verifiable
construction of the elimination hierarchy, a cost analysis tied to the
operations the solver executes, and an implementation that turns the analysis
into end-to-end speedups at high accuracy. The analysis uses only the temporal
locality of the relaxation, so it applies to multistage polynomial optimization
in general; in this paper we focus on robot trajectory problems, which motivate
the design and supply all of our test instances.

\paragraph{Organization.}
\Cref{sec:related} reviews related work. \Cref{sec:setting} fixes the SDP and
its Newton system. \Cref{sec:structure} derives temporal separators and an
explicit linear-work ordering, and \cref{sec:hierarchy} develops the certified
hierarchy. \Cref{sec:complexity} gives the complexity bounds, and
\cref{sec:algorithm} describes the GPU solver. \Cref{sec:experiments} reports
the experiments, and \cref{sec:discussion} discusses limitations. The appendices
treat factor reuse across iterations, the correspondence between the analysis and
the implementation, and additional measurements.
